\section{Experiments}
\label{sec:exp}

\subsection{Experimental Setup}

\myp{Datasets}
The statistics of the datasets used in the experiments are reported in Table~\ref{tab:datasets}.
\dsimdb is a KG describing actors and directors of movies.
\dsacm and \dsdblp are academic KGs denoting researchers and their publications with corresponding venues.
\dspubmed is a biomedical KG representing the relationships between genes, chemicals, and diseases of different species.
\dsfreebase is a general-purpose KG developed by Google.
\dsyelp is a KG built from user ratings and comments on businesses at different locations.
\dsbpm represents a series of synthetic datasets generated by the Bipartite Preferential Model~\cite{xirogiannopoulos2017extracting,guillaume2004bipartite,guillaume2006bipartite} for scalability tests (Sect.~\ref{sec:scalable}).

\input{tables/net_stats}

\vspace{1mm}
\myp{Query Generation}
For each dataset, we use \textsf{AnyBURL} \cite{meilicke2019anytime} to extract a set of candidate meta-paths $\mathbb{M}$.
We set the confidence threshold to $0.1$ for \textsf{AnyBURL} and take the snapshot at ten seconds.
Note that \textsf{AnyBURL} can extract meta-paths with extra node constraints, resulting in a large number of potential meta-paths for validation.
This makes efficient processing of \ced queries even more necessary. $|\mathbb{M}|$ in Table~\ref{tab:datasets} shows the number of meta-paths found on each dataset.
For personalized \ced queries, we randomly sample 10 nodes from $\vset_0$ as query nodes for each meta-path.

\vspace{1mm}
\myp{Compared Methods}
To the best of our knowledge, there have been no prior studies on the \ced problem.
Therefore, we compare our proposed methods with exact algorithms based on the state-of-the-art materialization methods.
The compared methods are listed as follows:
\begin{itemize}[leftmargin=*]
  \item \wcd and \wch find exact hub nodes by computing the respective degree and h-indexes of each node using the generic worst-case optimal join~\cite{ngo2018worst, freitag2020adopting, arroyuelo2021worst, veldhuizen2014leapfrog} (\S~\ref{sec:wcoj}).
  \item \boolap \cite{guo2024efficient} returns exact hubs based on both degree and h-index by materializing the relational graph through boolean matrix multiplication. {\sffamily BoolAP$^+$} is a variant of {\sffamily BoolAP} specially optimized for KGs with locally dense regions.
  \item \glod and \perd apply the sketch propagation framework (\S~\ref{sec:sketch-framework}) to estimate the degrees of all nodes in $H$. \glod returns the set of hubs based on estimated degrees, whereas \perd decides whether a query node $q$ is a hub.
  \item \perde optimizes \perd with early termination (\S~\ref{sec:degree-personal}).
  \item \gloh and \perh apply the pivot-based algorithm (\S~\ref{sec:hindex-global}) based on estimated h-indexes. \gloh recursively partitions the nodes until all hubs are found, while \perh checks if a query node $q$ is a hub using $q$ as a pivot.
  \item \perhe optimizes \perh with early termination (\S~\ref{sec:hindex-personal}).
\end{itemize}

\begin{table}[t]
    \setlength\tabcolsep{5pt}
    \footnotesize
    \centering
    \caption{Statistics of meta-paths in the experiments, where $|\eset^*|$ and $\rho^*$ denote the average number of edges and the density of matching graphs induced by meta-paths on each dataset.}\label{tab:EstarNumEdge}
    \vspace{-1em}
    \begin{tabular}{|c|c|c|c|c|c|}
        \hline
        & \dsimdb & \dsacm & \dsdblp & \dspubmed & \dsfreebase \\
        \hline
        $|\eset^*|$ & 22.3 & 133154  & 28941.4 & 23588.5 & 63526.3 \\
        \hline
        $\rho^*$ & 0.975 & 22.64  & 1.598 & 1.614 & 1.808 \\
        \hline
    \end{tabular}
    \vspace{-1em}
\end{table}

\input{tables/efficiency}

\myp{Parameter Settings}
By default, we set $\lambda = 0.05$ to find the hubs as the top 5\% of the nodes with the highest degrees or h-indexes in a relational graph.
Our results in \S~\ref{sec:effective} will show that h-index-based \ced queries are well approximated using smaller values of $\theta$ and $k$ compared to degree-based \ced queries.
Hence, for algorithms targeting degree-based \ced queries (\glod, \perd, \perde), we set $\theta=8$ and $k=32$ by default, while for algorithms targeting h-index-based \ced queries (\gloh, \perh, \perhe), we set $\theta=8$ and $k=4$.
These settings also demonstrate that our proposed methods produce accurate results in practice with significantly less stringency than those outlined in the worst-case analysis.

\vspace{1mm}
\myp{Evaluation Metrics}
We evaluate the quality of hubs each algorithm returns for global \ced queries (\glod and \gloh) with the F1-score.
The algorithms return a set $S$ of $\lambda \cdot |V|$ nodes.
However, there might be more than $\lambda \cdot |V|$ nodes satisfying the criteria of Problem~\ref{def:hub} due to the tied values with $c(v^\lambda)$.
Thus, we exclude nodes with centrality scores equal to $c(v^\lambda)$ in the recall calculation.
To assess the effectiveness of \perd, \perde, \perh, and \perhe for personalized \ced queries, we record their accuracy in deciding whether a node $q$ has a centrality score equal to or greater than $c(v^\lambda)$. We report the average F1 and accuracy scores across all meta-paths in Figs.~\ref{fig:effective-d} and~\ref{fig:effective-h}.
To evaluate the theoretical correctness of our proposed methods, we report the value of $\epsilon$ in Fig.~\ref{fig:epsilon} to measure the approximation ratio our methods achieved.
For efficiency evaluation, we report the average CPU time of each algorithm to process a \ced query.

\vspace{1mm}
\myp{Environment}
All experiments were carried out on a Linux server with AMD EPYC 7643 CPU and 256 GB memory running Ubuntu 20.04.
All algorithms were implemented in C++ with \texttt{-O3} optimization and executed in a single thread.
Our source code and data are published at \url{https://github.com/YudongNiu/HUB}.

\subsection{Efficiency Evaluation}
\label{sec:efficiency}

\begin{figure}[t]
    \centering
    \includegraphics[width=7.5cm]{figure/exp/efficiency-varylen-legend.pdf}\\
    \vspace{-1.5mm}
    \subfigure[Degree]{
    \begin{minipage}{4.1cm}
    \includegraphics[width=3.6cm]{figure/exp/efficiency-yelp-varylen-d.pdf}
    \end{minipage}
    \label{fig:yelp-d-varylen}
    }
    \subfigure[H-index]{
    \begin{minipage}{4.1cm}
    \includegraphics[width=3.6cm]{figure/exp/efficiency-yelp-varylen-h.pdf}
    \end{minipage}
    \label{fig:yelp-h-varylen}
    }
    \vspace{-1.5em}
    \caption{Running times of different algorithms on \dsyelp for meta-paths with varying length $L$.}
    \label{fig:varylen}
    \vspace{-2em}
\end{figure}

\begin{figure*}[t]
    \centering
    \includegraphics[width=4cm]{figure/exp/mpsensitivity/mp_sensitivity_legend_d.pdf}
    \vspace{-2mm}
    \\
    \subfigure[\dsimdb]{
    \begin{minipage}{3.1cm}
    \includegraphics[height=2.4cm]{figure/exp/mpsensitivity/imdb_mp_sensitivity_efficiency_df1.pdf}
    \end{minipage}
    \label{fig:imdb-d-mpsens}
    }
    \subfigure[\dsacm]{
    \begin{minipage}{3.1cm}
    \includegraphics[height=2.4cm]{figure/exp/mpsensitivity/acm_mp_sensitivity_efficiency_df1.pdf}
    \end{minipage}
    \label{fig:acm-d-mpsens}
    }
    \subfigure[\dsdblp]{
    \begin{minipage}{3.1cm}
    \includegraphics[height=2.4cm]{figure/exp/mpsensitivity/dblp_mp_sensitivity_efficiency_df1.pdf}
    \end{minipage}
    \label{fig:dblp-d-mpsens}
    }
    \subfigure[\dspubmed]{
    \begin{minipage}{3.1cm}
    \includegraphics[height=2.4cm]{figure/exp/mpsensitivity/pubmed_mp_sensitivity_efficiency_df1.pdf}
    \end{minipage}
    \label{fig:pubmed-d-mpsens}
    }
    \subfigure[\dsfreebase]{
    \begin{minipage}{3.1cm}
    \includegraphics[height=2.4cm]{figure/exp/mpsensitivity/freebase_mp_sensitivity_efficiency_df1.pdf}
    \end{minipage}
    \label{fig:freebase-d-mpsens}
    }
    \\
    \vspace{-1.5mm}
    \includegraphics[width=4cm]{figure/exp/mpsensitivity/mp_sensitivity_legend_h.pdf}\\
    \vspace{-2mm}
    \subfigure[\dsimdb]{
    \begin{minipage}{3.1cm}
    \includegraphics[height=2.4cm]{figure/exp/mpsensitivity/imdb_mp_sensitivity_efficiency_hf1.pdf}
    \end{minipage}
    \label{fig:imdb-h-mpsens}
    }
    \vspace{-1mm}
    \subfigure[\dsacm]{
    \begin{minipage}{3.1cm}
    \includegraphics[height=2.4cm]{figure/exp/mpsensitivity/acm_mp_sensitivity_efficiency_hf1.pdf}
    \end{minipage}
    \label{fig:acm-h-mpsens}
    }
    \subfigure[\dsdblp]{
    \begin{minipage}{3.1cm}
    \includegraphics[height=2.4cm]{figure/exp/mpsensitivity/dblp_mp_sensitivity_efficiency_hf1.pdf}
    \end{minipage}
    \label{fig:dblp-h-mpsens}
    }
    \subfigure[\dspubmed]{
    \begin{minipage}{3.1cm}
    \includegraphics[height=2.4cm]{figure/exp/mpsensitivity/pubmed_mp_sensitivity_efficiency_hf1.pdf}
    \end{minipage}
    \label{fig:pubmed-h-mpsens}
    }
    \subfigure[\dsfreebase]{
    \begin{minipage}{3.1cm}
    \includegraphics[height=2.4cm]{figure/exp/mpsensitivity/freebase_mp_sensitivity_efficiency_hf1.pdf}
    \end{minipage}
    \label{fig:freebase-h-mpsens}
    }
    \vspace{-0.5em}
    \caption{Running time of our methods for each meta-path on \dsfreebase with x-axis representing the number of edges $|\eset^*|$ in the corresponding matching graph. Lines of corresponding colors indicate the linear regression results of each method's running time with respect to $|\eset^*|$. Subfigures (a)-(e): degree based hub queries; Subfigures (f)-(j): h-index based hub queries.}
    \label{fig:mpsens}
    \vspace{-1.8em}
\end{figure*}

In Table~\ref{tab:efficiency}, we present the execution time per query for each method, along with its speedup ratios compared to the fastest exact method across five datasets.
Based on these results, we make the following observations.
First, sketch propagation-based algorithms are consistently more efficient than exact baselines \wcd and \wch across all datasets except \dsimdb and provide more than 10$\times$ speedups for degree-based queries (on \dsfreebase, \glod is 29$\times$ faster than \wcd) and up to two orders of magnitude speedups for h-index-based queries (on \dsfreebase, \gloh is 191$\times$ faster than \wch).
The speedup over degree-based queries is smaller than h-index-based queries because h-index-based queries are accurately approximated with smaller $k$ as discussed in \S~\ref{sec:effective}.
\glod is slightly slower than \wcd for degree-based queries on \dsimdb since matching graphs of meta-paths on \dsimdb are of relatively small sizes. Specifically, as demonstrated in Table~\ref{tab:EstarNumEdge}, the average number of edges in matching graphs on \dsimdb is 22.3, which is much smaller than matching graphs of meta-paths on other datasets. Therefore, the worst-case optimal join is performed quickly over the matching graph on \dsimdb, whereas our methods require additional costs for sketch construction and degree estimation.
Nevertheless, \glod answers \ced queries on \dsimdb in 1.3 milliseconds.
\boolapp, which is specially optimized for relational graph materialization based on matching graphs with dense regions, demonstrates exceptional efficiency on \dsacm, where the matching graphs have a high average density of 22.64.
Nevertheless, \boolapp is much slower than our sketch propagation methods across other datasets, where the matching graphs are of lower average density, and is less scalable even than \wcd and \wch.
Thus, in the following, we will only compare our methods with \wcd and \wch.
Second, the algorithms for personalized queries, i.e., \perde and \perhe, benefit from early termination optimization and achieve more than 2$\times$ additional speedups over \perd and \perh.

\myp{Effect of Selected Meta-Paths}
We investigate how the properties of specific meta-paths affect the efficiency of different methods.
Fig.~\ref{fig:varylen} presents the running time of different methods on \dsyelp varying the length of the meta-path $L$. We can observe that $L$ has a significant impact on efficiency, and exact methods \wcd and \wch exceed the 1-hour limit per meta-path when $L=2,3$. On the other hand, our methods based on \emph{sketch propagation} can efficiently process all queries up to $L=3$. Specifically, \glod and \gloh can answer \ced queries based on degree and h-index in 1,438 and 114 seconds, respectively. For personalized queries, \perhe and \perde handle meta-paths of $L=3$ in 893 and 42 seconds.

Fig.~\ref{fig:mpsens} presents the running time of our sketch propagation-based methods for each meta-path across datasets, along with the linear regression results of each method's running time with respect to the number of edges $|\eset^*|$ in the corresponding matching graph. We have the following observations. \underline{\textit{First},} the running time of all our methods generally increases with growing $|\eset^*|$, which aligns with the fact that the time complexity of our method is proportional to $|\eset^*|$.
\underline{\textit{Second},} the running time of \glod and \perd exhibits a strong correlation with $|\eset^*|$. For example, on \dsfreebase, we have $R^2$ coefficients of 0.91 and 0.78 for \glod and \perd, respectively. On the other hand, the running time of \perde demonstrates a weaker correlation with $|\eset^*|$, with a $R^2$ coefficient of 0.57 on \dsfreebase. By employing early termination during sketch propagation, \perde avoids propagating sketches over the entire matching graph, thereby reducing its dependence on $|\eset^*|$. A similar observation can be made for \gloh and \perh, whose running time is more closely correlated with $|\eset^*|$ compared to \perhe.
\underline{\textit{Third},} the running time of \glod and \perd generally exhibits a stronger correlation with $|\eset^*|$ compared to the running time of \gloh and \perh. For example, on \dspubmed, the $R^2$ coefficient is 0.91 for \glod and 0.86 for \gloh. This difference arises because \gloh requires multiple iterations of the sketch propagations based on randomly selected pivot nodes, making its running time dependent on both the number of sketch propagations and the time consumption for each propagation, which is influenced by $|\eset^*|$. In contrast, \glod requires a single pass of sketch propagation for the KMV sketch construction.

\leavetotechreport{Due to space limit, the results for memory consumption are deferred to the technical report~\cite{hud2023tr}.}
We exclude \dsyelp from the remaining experiments on effectiveness evaluations to ensure a fair comparison, as \wcd and \wch cannot complete on all meta-paths.

\techreport{We then report the memory consumption of different methods in Fig.~\ref{fig:mem}. We can observe from Fig.~\ref{fig:mem-d} that \glod, \perd, and \perde take relatively more memory than the exact algorithm \wcd to maintain KMV sketches for images. Nevertheless, our methods take at most 1.93 times more memory than \wcd (on \dsdblp). On the other hand, the gap between the memory consumption of \gloh, \perh, \perhe and the memory consumption of \wch is marginal since h-index-based queries require fewer KMV sketches for accurate estimation. Specifically, our methods take at most 1.21 times more memory than \wch (on \dspubmed).

\begin{figure}[t]
    \centering
    \includegraphics[width=8cm]{figure/exp/efficiency-varylen-legend.pdf}
    \subfigure[degree]{
        \begin{minipage}{3.8cm}
        \includegraphics[width=4cm]{figure/exp/mem-d.pdf}
        \end{minipage}
        \label{fig:mem-d}
    }
    \subfigure[h-index]{
        \begin{minipage}{3.8cm}
        \includegraphics[width=4cm]{figure/exp/mem-h.pdf}
        \end{minipage}
        \label{fig:mem-h}
    }
    \caption{Memory consumption of compared methods.}
    \label{fig:mem}
\end{figure}
}

\begin{figure*}[t]
    \centering
    \includegraphics[width=4.5cm]{figure/exp/effect-global/legend-d.pdf}
    \vspace{-1mm}
    \\
    \subfigure[\dsimdb]{
	\begin{minipage}{3.1cm}
        \includegraphics[height=2.45cm]{figure/exp/effect-global/effect-d-imdb-varying-R.pdf}
        \end{minipage}
        \label{fig:imdb-d-varyingR}
    }
    \vspace{-3mm}
    \subfigure[\dsacm]{
        \begin{minipage}{3.1cm}
        \includegraphics[height=2.45cm]{figure/exp/effect-global/effect-d-acm-varying-R.pdf}
        \end{minipage}
        \label{fig:acm-d-varyingR}
    }
    \subfigure[\dsdblp]{
	\begin{minipage}{3.1cm}
        \includegraphics[height=2.45cm]{figure/exp/effect-global/effect-d-dblp-varying-R.pdf}
        \end{minipage}
        \label{fig:dblp-d-varyingR}
    }
    \subfigure[\dspubmed]{
        \begin{minipage}{3.1cm}
        \includegraphics[height=2.45cm]{figure/exp/effect-global/effect-d-pubmed-varying-R.pdf}
        \end{minipage}
        \label{fig:pubmed-d-varyingR}
    }
    \subfigure[\dsfreebase]{
        \begin{minipage}{3.1cm}
        \includegraphics[height=2.45cm]{figure/exp/effect-global/effect-d-freebase-varying-R.pdf}
        \end{minipage}
        \label{fig:freebase-d-varyingR}
    }
    \subfigure[\dsimdb]{
        \begin{minipage}{3.1cm}
        \includegraphics[height=2.45cm]{figure/exp/effect-global/effect-d-imdb-varying-K.pdf}
        \end{minipage}
        \label{fig:imdb-d-varyingk}
    }
    \vspace{-1mm}
    \subfigure[\dsacm]{
        \begin{minipage}{3.1cm}
        \includegraphics[height=2.45cm]{figure/exp/effect-global/effect-d-acm-varying-K.pdf}
        \end{minipage}
        \label{fig:acm-d-varyingk}
    }
    \subfigure[\dsdblp]{
	\begin{minipage}{3.1cm}
        \includegraphics[height=2.45cm]{figure/exp/effect-global/effect-d-dblp-varying-K.pdf}
        \end{minipage}
        \label{fig:dblp-d-varyingk}
    }
    \subfigure[\dspubmed]{
        \begin{minipage}{3.1cm}
        \includegraphics[height=2.45cm]{figure/exp/effect-global/effect-d-pubmed-varying-K.pdf}
        \end{minipage}
        \label{fig:pubmed-d-varyingk}
    }
    \subfigure[\dsfreebase]{
        \begin{minipage}{3.1cm}
        \includegraphics[height=2.45cm]{figure/exp/effect-global/effect-d-freebase-varying-K.pdf}
        \end{minipage}
        \label{fig:freebase-d-varyingk}
    }
    \vspace{-0.5em}
    \caption{Effectiveness of \glod, \perd, and \perde for \ced queries based on degree centrality with varying parameters.
    Subfigures (a)-(e): varying parameter $\lambda$. Subfigures (f)-(j): varying parameter $k$.} % Subfigures (k)-(o): varying parameter $\theta$.}
    \label{fig:effective-d}
    \vspace{-1em}
\end{figure*}

\begin{figure*}[t]
    \centering
    \includegraphics[width=4.5cm]{figure/exp/effect-global/legend-h.pdf}
    \vspace{-1mm}
    \\
    \subfigure[\dsimdb]{
    \begin{minipage}{3.1cm}
    \includegraphics[height=2.45cm]{figure/exp/effect-global/effect-h-imdb-varying-R.pdf}
    \end{minipage}
    \label{fig:imdb-h-varyingR}
    }
    \vspace{-3mm}
    %
    \subfigure[\dsacm]{
    \begin{minipage}{3.1cm}
    \includegraphics[height=2.45cm]{figure/exp/effect-global/effect-h-acm-varying-R.pdf}
    \end{minipage}
    \label{fig:acm-h-varyingR}
    }
    \subfigure[\dsdblp]{
    \begin{minipage}{3.1cm}
    \includegraphics[height=2.45cm]{figure/exp/effect-global/effect-h-dblp-varying-R.pdf}
    \end{minipage}
    \label{fig:dblp-h-varyingR}
    }
    \subfigure[\dspubmed]{
    \begin{minipage}{3.1cm}
    \includegraphics[height=2.45cm]{figure/exp/effect-global/effect-h-pubmed-varying-R.pdf}
    \end{minipage}
    \label{fig:pubmed-h-varyingR}
    }
    \subfigure[\dsfreebase]{
    \begin{minipage}{3.1cm}
    \includegraphics[height=2.45cm]{figure/exp/effect-global/effect-h-freebase-varying-R.pdf}
    \end{minipage}
    \label{fig:freebase-h-varyingR}
    }
    \subfigure[\dsimdb]{
    \begin{minipage}{3.1cm}
    \includegraphics[height=2.45cm]{figure/exp/effect-global/effect-h-imdb-varying-K.pdf}
    \end{minipage}
    \label{fig:imdb-h-varyingk}
    }
    \vspace{-1mm}
    \subfigure[\dsacm]{
    \begin{minipage}{3.1cm}
    \includegraphics[height=2.45cm]{figure/exp/effect-global/effect-h-acm-varying-K.pdf}
    \end{minipage}
    \label{fig:acm-h-varyingk}
    }
    \subfigure[\dsdblp]{
    \begin{minipage}{3.1cm}
    \includegraphics[height=2.45cm]{figure/exp/effect-global/effect-h-dblp-varying-K.pdf}
    \end{minipage}
    \label{fig:dblp-h-varyingk}
    }
    \subfigure[\dspubmed]{
    \begin{minipage}{3.1cm}
    \includegraphics[height=2.45cm]{figure/exp/effect-global/effect-h-pubmed-varying-K.pdf}
    \end{minipage}
    \label{fig:pubmed-h-varyingk}
    }
    \subfigure[\dsfreebase]{
    \begin{minipage}{3.1cm}
    \includegraphics[height=2.45cm]{figure/exp/effect-global/effect-h-freebase-varying-K.pdf}
    \end{minipage}
    \label{fig:freebase-h-varyingk}
    }
    \vspace{-0.5em}
    \caption{Effectiveness of \gloh, \perh, and \perhe for \ced queries based on h-index with varying parameters. Subfigures (a)-(e): varying parameter $\lambda$. Subfigures (f)-(j): varying parameter $k$.} % Subfigures (k)-(o): varying parameter $\theta$.}
    \label{fig:effective-h}
    \vspace{-1em}
\end{figure*}

\begin{figure*}[t]
    \centering
    \includegraphics[width=4.5cm]{figure/exp/epsilon/epsilon_legend.pdf}
    \vspace{-2mm}
    \\
    \subfigure[\dsimdb]{
    \begin{minipage}{3.1cm}
    \includegraphics[height=2.45cm]{figure/exp/epsilon/epsilon_imdb.pdf}
    \end{minipage}
    \label{fig:imdb-epsilon}
    }
    %
    \subfigure[\dsacm]{
    \begin{minipage}{3.1cm}
    \includegraphics[height=2.45cm]{figure/exp/epsilon/epsilon_acm.pdf}
    \end{minipage}
    \label{fig:acm-epsilon}
    }
    \subfigure[\dsdblp]{
    \begin{minipage}{3.1cm}
    \includegraphics[height=2.45cm]{figure/exp/epsilon/epsilon_dblp.pdf}
    \end{minipage}
    \label{fig:dblp-epsilon}
    }
    \subfigure[\dspubmed]{
    \begin{minipage}{3.1cm}
    \includegraphics[height=2.45cm]{figure/exp/epsilon/epsilon_pubmed.pdf}
    \end{minipage}
    \label{fig:pubmed-epsilon}
    }
    \subfigure[\dsfreebase]{
    \begin{minipage}{3.1cm}
    \includegraphics[height=2.45cm]{figure/exp/epsilon/epsilon_freebase.pdf}
    \end{minipage}
    \label{fig:freebase-epsilon}
    }
    \vspace{-0.5em}
    \caption{Average approximation ratio $\epsilon$ achieved by our methods by varying the sketch size $k$, with shaded area denoting the variance of $\epsilon$ over 10 executions.}
    \label{fig:epsilon}
    \vspace{-1em}
\end{figure*}

\subsection{Effectiveness Evaluation}
\label{sec:effective}

\myp{Results for Degree-based \ced Queries}
Figs.~\ref{fig:imdb-d-varyingR}--\ref{fig:freebase-d-varyingR} illustrate the effectiveness of \glod, \perd, and \perde on each dataset by varying the parameter $\lambda$.
We observe that \glod achieves F1-scores of more than 0.85 across all datasets for \ced queries, and \perd and \perde consistently achieve accuracy rates of at least 0.98 for personalized \ced queries.
Moreover, \glod provides more accurate answers with increasing $\lambda$ and achieves an F1-score of over 0.9 when $\lambda=0.05$.
The reason is that for a larger $\lambda$, the $(1-\lambda)$-quantile node $v^\lambda$ has a smaller degree, and \ced queries apply less strict requirements to the hubs, which can be more easily approximated. The effectiveness of \glod varies across datasets: It achieves higher F1-scores on \dsimdb and \dsacm than on the other three datasets. The variation in the effectiveness of \glod is caused by different degree distributions of relational graphs in each dataset. Specifically, matching graphs with higher density on \dsacm result in densely connected relational graphs, where on average 72.7\% of the nodes share the same degree as $v^\lambda$. These nodes are excluded from the recall calculation, leading to a higher F1-score. In contrast, the percentage of nodes with the same degree as $v^\lambda$ is 32.25\% on \dsdblp, 38.64\% on \dspubmed, and 30.6\% on \dsfreebase. \glod achieves a high F1-score on \dsimdb because the matching graphs are relatively small, allowing accurate degree estimation in the relational graphs.

Figs.~\ref{fig:imdb-d-varyingk}--\ref{fig:freebase-d-varyingk} show the effectiveness of \glod, \perd, and \perde on each dataset with varying the parameter $k$. 
First, \glod, \perd, and \perde all provide better results with increasing $k$. Larger KMV sketches allow the sketch propagation-based methods to approximate \ced queries with higher F1-scores or accuracy rates due to more accurate degree estimations.
Second, personalized \ced queries are accurately approximated by \perd and \perde, even with relatively smaller values of $k$.
This is because personalized \ced queries only require comparing the degrees of $q$ and $v^\lambda$, which is relatively easy due to the significant difference between the degrees of $q$ and $v^\lambda$.
Third, the gap between the accuracy of \perde and \perd is marginal, even when $k=2$ (with a maximum gap of 0.044 on \dsdblp), and narrows further with increasing $k$. This is because a larger $k$ leads to more accurate estimations of image sizes, thus reducing the chance of false early termination.

\vspace{1mm}
\myp{Results for \ced Queries based on H-index}
The effectiveness results of \gloh, \perh, and \perhe are presented in Fig.~\ref{fig:effective-h}.
Similar observations can be made for the methods for h-index-based \ced queries as those for degree-based \ced queries.
Furthermore, we note that \gloh achieves a higher F1-score for h-index-based \ced queries compared to \glod in the same parameter setting.
This is because the range of h-index values is much smaller than that of degrees, resulting in a larger number of tied h-index values.
By excluding nodes with h-index values equal to $c(v^\lambda)$ from the recall calculation, h-index-based \ced queries become much easier to approximate.

% To evaluate the theoretical correctness of our methods in terms of the concept of $\epsilon$-separable, we further study the $\epsilon$ our methods achieved. 
\myp{Results for Approximation Ratio $\epsilon$}
Fig.~\ref{fig:epsilon} presents the approximation ratio $\epsilon$ achieved by our sketch propagation-based methods varying sketch size $k$. We have the following observations. \underline{\textit{First},} our methods answer Problems \ref{problem:approxHUB} and \ref{problem:approxPer} accurately with $\epsilon$ consistently remaining below 0.6 under the default parameter setting of $k$ and $\theta$. Specifically, the approximation ratio $\epsilon$ achieved by our methods decreases monotonically with the increase of $k$ across datasets and approaches zero when $k=32$ except on \dsacm. This exception arises due to relational graphs with a substantial proportion of high-degree nodes on \dsacm induced by matching graphs with high density, thereby requiring a larger value of $k$ to effectively identify hubs in these relational graphs, as indicated by Theorems~\ref{theorem:dglobal} and~\ref{theorem:hindex}. \underline{\textit{Second},} given fixed value of $k$, \gloh, \perh and \perhe can achieve smaller $\epsilon$ compared to \glod, \perd and \perde. This observation complies with our observation from Figs.~\ref{fig:effective-d} and~\ref{fig:effective-h} that $h$-index-based queries are easier to approximate due to the large number of tied $h$-index values in relational graphs. Note that \perd and \perh achieve the same approximation ratio as \glod and \gloh, respectively, since the query node for personalized hub queries can be selected as the node leading to the maximum $\epsilon$ in global hub queries.

\vspace{1mm}
\myp{Summary}
\glod, \perd, and \perde achieve F1-scores or accuracy rates greater than 0.9 across all datasets when $k=32$ and $\theta=8$ for degree-based \ced queries with $\lambda=0.05$.
Similarly, for h-index-based \ced queries, \gloh, \perh, and \perhe achieve F1 scores or an accuracy greater than 0.9 across all datasets when $k=4$ and $\theta=8$.
All of these results confirm the empirical effectiveness as well as theoretical correctness of our proposed methods.

\begin{figure}[t]
    \centering
    \includegraphics[width=5cm]{figure/exp/influence_analysis_legend.pdf}
    \\
    \includegraphics[width=7.5cm]{figure/exp/influence_analysis.pdf}
    \vspace{-4mm}
    \caption{Ratio between the influence scores achieved by hubs revealed by \ced methods and the highest influence scores estimated by the RIS algorithm.}
    \label{fig:influence}
    \vspace{-1em}
\end{figure}

\subsection{Influence Analysis}

We further verify the effectiveness of \ced for influence analysis on relational graphs.
We adopt the weighted cascade model for influence estimation~\cite{goyal2011data,chen2010scalable}.
Fig.~\ref{fig:influence} reports the ratio between the influence scores of top-5\% hubs returned by \ced methods (\wcd, \glod, \wch, and \gloh) and the influence scores of top-5\% seeds with the highest influence estimated by a reverse influence sampling (RIS) algorithm for influence maximization~\cite{tang2015influence}. We have the following observations.
First, all \ced methods consistently achieve more than 90\% of the influence scores achieved by the RIS algorithm, which verifies that the hubs have significant influence over the relational graph and that our proposed methods can serve as effective seeding strategies.
Second, the h-index-based hubs (obtained by \gloh and \wch) outperform the degree-based hubs (obtained by \glod and \wcd) and maintain over 95\% of the influence scores achieved by the RIS algorithm. In contrast to degree-based hubs, h-index-based hubs have neighbors with higher degrees, which further facilitate the spread of influence to indirect neighbors.

\subsection{Scalability Test}
\label{sec:scalable}

We further test the scalability of our methods on large synthetic datasets \dsbpm. Fig.~\ref{fig:bpm} reports the running time and memory consumption of different methods by varying the number of nodes. We make the following observations.
The exact methods \wcd and \wch cannot scale to large datasets. Specifically, \wcd and \wch can only process queries from the smallest \dsbpm datasets with less than 1M nodes and exceed the 1-hour limit on larger datasets. In contrast, our methods based on \emph{sketch propagation} scale linearly with the size of the datasets and can handle queries on \dsbpm datasets with up to 192M nodes and 1.92B edges within 2,552 seconds for degree-based queries and 258 seconds for h-index-based queries.
In terms of memory consumption, our methods use 118.5 GB and 134 GB memory, respectively, to process queries based on the h-index and degree centrality measures on the largest \dsbpm dataset. The difference between the memory consumption of queries based on different centrality measures is marginal because KMV sketches are memory efficient compared to the matching graph itself.

\begin{figure}[t]
    \centering
    \includegraphics[width=8cm]{figure/exp/bipartite-legend.pdf}\\
    \vspace{-1.5mm}
    \subfigure[Running Time]{
        \includegraphics[width=4.1cm]{figure/exp/bipartite-efficiency.pdf}
        \label{fig:scale:efficiency}
    }
    \subfigure[Memory Usage]{
        \includegraphics[width=4.1cm]{figure/exp/bipartite-mem.pdf}
        \label{fig:scale:mem}
    }
    \vspace{-1em}
    \caption{Runtime and memory consumption of our methods on the \dsbpm dataset.}
    \label{fig:bpm}
    \vspace{-5mm}
\end{figure}
